% Tutorial concepts consolidation — what a sheet taught, built from the area's records rather
% than from the solutions a second time. One per sheet, beside the writeup.
% Copy it there and replace every <...> placeholder. The input line below resolves the preamble
% from the templates folder or from a unit folder two levels down, so it never needs editing.
\documentclass[a4paper,10pt]{article}
\IfFileExists{preamble.tex}{\input{preamble.tex}}{\input{../../templates/preamble.tex}}
\begin{document}
\ArtifactHeader{\ModuleCode}{<Source-map key>}{<the sheet's own name> --- Concepts}{<YYYY-MM-DD>}%
  {<the sheet, and the records this was built from, by number>}

\section{What the sheet was for}
<The one thing the sheet was built to make land, in a sentence. If that sentence is hard to write,
the records are the place to go back to, not the solutions.>

\section{<The idea the sheet exercises>}

\begin{defn}[<term>]
<Restated in the form the questions actually used it, which is often narrower than the lecture's.>
\end{defn}

\begin{remark}
<The move the sheet kept asking for, named so it is recognisable next time it is needed.>
\end{remark}

\begin{warning}
<Where the working went wrong, taken from the records rather than guessed — this is the part a
later session reads first.>
\end{warning}

\section{Where this shows up again}
\begin{itemize}
  \item <A question elsewhere in the module that turns on the same idea, named by pointer.>
  \item <The unit whose material this sheet was exercising, by Source-map key.>
\end{itemize}

\begin{checkyourself}
  \item <The question from the sheet worth being able to do cold.>
\end{checkyourself}
\end{document}
